\documentclass[10pt,a4paper]{amsart}
\usepackage[margin=19mm]{geometry}
\usepackage{amsmath,amssymb,mathtools}
\usepackage[scr=rsfs,cal=euler]{mathalfa}
\usepackage{xfrac}
\usepackage[unicode,hidelinks,bookmarks=false]{hyperref}
\newcommand{\ie}{i.e.\ }
\newcommand{\cf}{cf.\ }
\newcommand{\resp}{respectively\ }
\newcommand{\Subsection}{\subsection}
\providecommand{\nobreakdash}{\nobreak\hbox{-}}
\setlength{\emergencystretch}{2em}
\multlinegap0pt
\usepackage{bm}
\usepackage{bbm}
\usepackage{dsfont}
\usepackage{xspace}
\usepackage{cancel}
\usepackage{mathtools}
\usepackage{amsthm}
\makeatletter
\@ifundefined{thm}{\newtheorem{thm}{\bf Theorem}}{}
\makeatother
\title[The Classification of Rotationally symmetric hypersurfaces i]{The Classification of Rotationally symmetric hypersurfaces in the Heisenberg groups $H\_{n}$}
\author{Hung-Lin Chiu, Sin-Hua Lai, Hsiao-Fan Liu}
\date{Source: arXiv:2509.04972v1 (CC BY 4.0)}
\begin{document}
\maketitle

\noindent\emph{Source and licence.} The Classification of Rotationally symmetric hypersurfaces in the Heisenberg groups $H_{n}$. Version arXiv:2509.04972v1 (CC BY 4.0), distributed under the Creative Commons Attribution 4.0 International licence, as recorded in the arXiv licence metadata for this exact version and shown on the abstract page https://arxiv.org/abs/2509.04972v1. Full rights notice: licenses/arxiv-2509.04972-CC-BY-4.0.txt.\par\smallskip

\noindent\textit{Editorial scope.} Counted unit: the Thm whose source label is \texttt{funthm1} in the article (source0.tex lines 494--497, source label \texttt{funthm1}), delivered together with the article's own complete original proof of it (source0.tex lines 498--541, one \texttt{proof} environment of 44 lines). No mathematics is written, repaired or re-derived: statement and proof are the article's own text line for line. Required context retained uncounted: regecu (lines 474--479); regecu1 (lines 482--487); nomal (lines 489--491). Every definition, display and label of those lines is retained unchanged. Omitted from this package and not redistributed: 1--473, 480--481, 488--488, 492--493, 542--857. Nothing inside a retained excerpt is omitted or altered: every retained body is delivered exactly as the article's lines are written, so no omission receipt of any kind accompanies this package. Citations: the retained excerpts carry no citation command at all, so no bibliography entry of the article is transcribed and no reference is redistributed. Reference adaptation: none was needed and none was made; every reference inside the delivered unit resolves inside it, so no unresolved reference or citation is produced. Printed slips kept literal: no printed word, symbol, hypothesis or number of the counted statement or of its proof was altered. Layout and class: the article's own class is \texttt{amsart} with packages including amssymb, amsmath, amscd, graphicx, mathrsfs, pgf, tikz; the wrapper is the lane's pinned \texttt{amsart} editorial wrapper at 10pt on A4, which restates the theorem-like environments the retained text needs and adds the article's own macro definitions that the retained text needs (none), with extra wrapper packages (bm, bbm, dsfont, xspace, cancel, mathtools). The delivered numbering is the unit's own. Assets: no retained span contains a figure environment, a graphics-inclusion command, a PDF-page inclusion, a table or a TikZ picture; no image file is copied into this package, the package declares no asset dependency and no image file is redistributed with it. Licence: version arXiv:2509.04972v1 of this article is distributed under the Creative Commons Attribution 4.0 International licence, as recorded in the arXiv licence metadata for this exact version and shown on the abstract page https://arxiv.org/abs/2509.04972v1. A full rights notice is delivered at licenses/arxiv-2509.04972-CC-BY-4.0.txt. Characters: every retained excerpt is pure ASCII, so the delivered engine, XeLaTeX with its default Latin Modern text font, reports no missing character.
\begin{equation}\label{regecu}
\begin{split}
x(\tilde{s})&=\int \alpha x(\tilde{s})d\tilde{s}\ \ \left(\Leftrightarrow x(\tilde{s})=e^{\int\alpha d\tilde{s}}\right)\\
t(\tilde{s})&=-\int kx^2(\tilde{s})d\tilde{s},
\end{split}
\end{equation}

\begin{equation}\label{regecu1}
\begin{split}
x(s)&=e^{\int\bar\alpha ds},\\
t(s)&=-\int \bar{k}x^2 ds,
\end{split}
\end{equation}

\begin{equation}\label{nomal}
\frac{(xx')^{2}+(t')^2}{x^2}=1.
\end{equation}

\begin{thm}\label{funthm1}
Given two functions $\bar\alpha(s),\bar{k}(s)$ which satisfy the integrability condition \eqref{nomal}.
Then there exists a rotationally symmetric hypersurface such that $\alpha(s)=\bar\alpha(s)$ and $k(s)=\bar k(s)$ for all $s$. And $s$ is the horizontal arc-length. In addition, such a hypersurface is unique, up to a Heisenberg rigid motion.
\end{thm}

\begin{proof}
We consider the horizontal generating curve 
\[\tilde{\gamma}(s)=(0,\cdots,0,x(s)\cos{\theta(s)},0,\cdots,0,x(s)\sin{\theta(s)},t(s)).\]
for some angle function $\theta=\theta(s)$. After a straightforward computation, we have
\[\begin{split}
\frac{d\tilde{\gamma}}{ds}(s)&=(x'\cos{\theta}-x\sin{\theta}\cdot\theta')\frac{\partial}{\partial x_n}+(x'\sin{\theta}+x\cos{\theta}\cdot\theta')\frac{\partial}{\partial y_n}+(t')\frac{\partial}{\partial t}\\
&=(x'\cos{\theta}-x\sin{\theta}\cdot\theta')\mathring{e}_{n}+(x'\sin{\theta}+x\cos{\theta}\cdot\theta')\mathring{e}_{2n}+(x^2\theta'+t')\frac{\partial}{\partial t},
\end{split}\]
which implies that $\tilde{\gamma}(s)$ is horizontal if and only if 
\begin{equation*}
\theta'=-\frac{t'}{x^2}.
\end{equation*}
Thus, we have
\[\begin{split}
\left\|\frac{d\tilde{\gamma}}{ds}\right\|^{2}&=(x'\cos{\theta}-x\sin{\theta}\cdot\theta')^{2}+(x'\sin{\theta}+x\cos{\theta}\cdot\theta')^{2}\\
&=(x')^{2}+x^{2}(\theta')^{2}=\frac{(xx')^2+(t')^2}{x^2},
\end{split}\]
which shows that the normalization \eqref{nomal} just means that $s$ is the horizontal arc-length $\tilde{s} $ of the corresponding horizontal generating curve. Next, from the definition of the generating curve \eqref{regecu1}, we immediately have 
\begin{equation}\label{goinf4}
\bar\alpha=x^{-1}x'\ \ \textrm{and}\ \  \bar k=-x^{-2}t'. 
\end{equation}
Thus,
\[\begin{split}
\bar\alpha^{2}+\bar k^{2}&=x^{-2}(x')^{2}+x^{-4}(t')^{2}\\
&=\frac{1}{x^2}\left(\frac{(xx')^{2}+(t')^{2}}{x^{2}}\right)=\frac{1}{x^2}.
\end{split}\]
Let $\alpha$ and $k$ be the $\alpha$-function and $k$-function of the generated surface. Then \eqref{regecu} and \eqref{regecu1} yield that
\[(\alpha^2+k^{2})(\tilde{s})=(\bar\alpha^2+\bar k^{2})(s)\] and
\begin{equation}\label{regecu2}
\begin{split}
t(\tilde{s})&=-\int kx^2(\tilde{s})d\tilde{s}=-\int \bar kx^2(s)ds\\
&=-\int \bar kx^2(s)\frac{ds}{d\tilde{s}}d\tilde{s},
\end{split}
\end{equation}
for all $\tilde{s}$.
Therefore, we have
\begin{equation}
\begin{split}
k(\tilde{s})&=\bar k(s)\frac{ds}{d\tilde{s}},\\
\alpha^2(\tilde{s})&=\bar\alpha^2(s)+\bar k^{2}\left(1-(\frac{ds}{d\tilde{s}})^2\right).
\end{split}
\end{equation}
In particular, we have $k(\tilde{s})=\bar k(s)$ and $\alpha^2(\tilde{s})=\bar\alpha^2(s)$ as $\frac{ds}{d\tilde{s}}=1$, i.e, $s$ is the horizontal arc-length of the corresponding horizontal generating curve. Since \[x(s)=e^{\int \alpha ds}=e^{\int\bar\alpha ds},\] we have that $\alpha(\tilde{s})=\bar\alpha(s)$. This completes the proof of Theorem \ref{funthm1} and thus provides a proof for {\bf Theorem I}.
\end{proof}

\end{document}
